\begin{abstract}
Failure predictors for shared GPU clusters are often evaluated on random task-level splits. We audit that practice on the Alibaba PAI 2020 trace with KANVAS, an additive Kolmogorov-Arnold model (\kaam{}) whose ten learned B-spline curves are its complete explanation. Ten leak-free features are built for every labelled task by averaging other jobs' telemetry on the task's machine in the hour before it started; \V{data.full.tasks} tasks (\V{data.full.share}) pass the filters. Clipping and scaling are fitted on training splits only, which removes a leak in our earlier version. Under four protocols fixed before any test label was read, we compare \kaam{} with logistic regression, a parameter-matched MLP, LightGBM and, on an earlier \V{data.orig.tasks}-task sample, two Explainable Boosting Machines, using paired bootstrap tests with Holm correction. On random splits LightGBM leads (\V{full.random.lgbm.auc} ROC-AUC against \V{full.random.kaam.auc} for \kaam{}), but \V{rf.full.dup.random.min} of the random-split test tasks have an exact copy of all ten features in training. With users held out the lead disappears: on the four user-grouped folds without the dominant user, the MLP (\V{full.gnou.mlp.auc}) and \kaam{} (\V{full.gnou.kaam.auc}) score above LightGBM (\V{full.gnou.lgbm.auc}), and no model is clearly best. On a time-ordered split every model scores between \V{full.time.minauc} and \V{full.time.maxauc}, even on users seen in training. \kaam{}'s curves and exact per-task contributions expose a shortcut: one-core requests carry a large risk term because one user, who submits \V{users.full.u1.taskshare} of tasks and \V{users.full.u1.failshare} of failures, holds most one-core training tasks; without that user, one-core tasks fail about as often as others. The features rank tasks well only for users and periods seen in training, and random-split accuracy overstates performance on new users and on a later period.
\end{abstract}
